% ======================================================================
% SECTION 1: HR 9501 - Cancer Patient Self-Selection of Physical AI
% Oncology Trials Act of 2026
% Adaption of 21 CFR Part 50, FDA Decentralized Trial Final Guidance
% 2024, and 42 USC 300gg-8.
%
% Final prose populated from the consolidated bracketed instructions
% in the prior commit. The five-subsection legislative-act layout
% follows the SB 1042 pattern from
% new-trial/site/all-documents/all_documents.tex.
% ======================================================================

\hspace{0.5cm} The cancer patient, not the doctor, the nurse, the trial sponsor, the
IRB, or the regulator, is the priority participant in any United
States oncology clinical trial. Advanced AI and advanced robotics now
offer the patient direct control over discovery, enrollment,
scheduling, robot selection, procedural modification, real-time data
feedback, and self-custody of health data. HR 9501 is the foundational
bill of seven proposed federal bills (HR 9501 through HR 9507) that
supply the legal infrastructure for that control. The bill is an
\textit{Adaption of 21 CFR Part 50, FDA Decentralized Clinical Trial
Final Guidance 2024, and 42 USC 300gg-8}, drawing the chain of prior
authority from the Common Rule baseline through the FDA decentralized
elements guidance and the Medicaid coverage statute that already binds
qualifying clinical trials \cite{hhs_45cfr46_protection_human_subjects,
fda_dct_guidance_2024, fda_oce_advancing_oncology_dct_2024,
usc_300gg8_clinical_trials, kawchak_2026_19994945,
kawchak_2026_19119939, kawchak_2026_20045457}.

The current legislative stack supplies the points of departure for
this bill. The CLINICAL TREATMENT Act of 2020 (PL 116-260) established
Medicaid coverage of routine costs in qualifying trials from January
1, 2022; HR 9501 extends that coverage right into a positive
self-selection right that includes 24/7 booking, decentralized payment
parity, and AI trial-matching tool actionability for Medicaid-eligible
patients \cite{rank2_clinical_treatment_act_2020, jin2024trialgpt,
gupta2024prism, abdallah2026trialmatchai}. Affordable Care Act Section
2709 (PL 111-148, 2010) set the private-insurance baseline for routine
clinical-trial cost coverage; HR 9501 closes the decentralized and
at-home component payment-parity gap that ACA Section 2709 left
unaddressed \cite{rank3_aca_clinical_trials_2010,
usc_300gg8_clinical_trials}. FDAAA Section 801 (PL 110-85, 2007)
established the ClinicalTrials.gov infrastructure; HR 9501, jointly
with HR 9507, closes the structured-eligibility-data gap that prevents
patient-initiated trial discovery \cite{rank5_fdaaa_2007}. HHS HIPAA
Right-to-Access (2025), the ASTP/ONC Patient Portals Brief 2024, and
the ASTP/ONC Hospital APIs Brief 2024 establish the foundational
data-access framework on which self-selection depends; the self-
custody right itself is reserved for HR 9507
\cite{hhs_right_to_access_2025, astp_onc_patient_portals_apps_2024,
astp_onc_hospital_apis_2024}.

HR 9501 implements two of the seven dimensions of the operational
definition of patient control that anchor this paper. The trial-discovery dimension (machine-readable eligibility, AI matching, no
provider gatekeeping) is implemented through the Subsection 1.3 rights
below; the location and scheduling dimension (24/7 booking,
decentralized elements, home-based delivery) is implemented through
the same operative core \cite{jin2024trialgpt, gupta2024prism,
abdallah2026trialmatchai, fda_dct_guidance_2024,
fda_oce_advancing_oncology_dct_2024, dronca2026cancer_care_beyond_walls,
kawchak_2026_19994945}.

\subsection{Findings and Declarations}

\hspace{0.5cm} Congress finds and declares the following.

(a) Approximately seven percent of adults with cancer in the United
States enroll in an oncology clinical trial, despite documented
patient-benefit advantages of trial participation
\cite{kawchak_2026_19994945}. The participation gap arises in part
because the prior framework requires provider-initiated trial referral;
patient-initiated discovery is technically possible but is not yet a
statutory right \cite{rank2_clinical_treatment_act_2020,
rank3_aca_clinical_trials_2010}.

(b) Large language model-based trial-matching systems have demonstrated
recall and accuracy at or above the clinician baseline. TrialGPT, a
Nature Communications 2024 system, achieved a top-10 recall of 87.3
percent on a held-out test set; PRISM extended these numbers across
multi-site EHR data; TrialMatchAI was published in Nature
Communications 2026 with end-to-end recommendation latency on the
order of seconds \cite{jin2024trialgpt, gupta2024prism,
abdallah2026trialmatchai}.

(c) The FDA Decentralized Clinical Trials Final Guidance 2024, the
FDA Oncology Center of Excellence Advancing Oncology Decentralized
Trials directory, and the FDA Real-Time Clinical Trials press
announcement of April 28, 2026 together establish a regulatory
opening for patient-initiated trial selection at sites that admit
decentralized elements \cite{fda_dct_guidance_2024,
fda_oce_advancing_oncology_dct_2024, fda_rtct_press_2026}.

(d) The patient-side technical substrate is documented in the
author's prior 24/7 continuous trial paper at
\texttt{new-trial/national-24-7-trial/paper/full-paper/}, in the
single-patient 10-stage journey at \texttt{patient-journey/}, and in
the autonomous sponsor 168-hour extension at
\texttt{sponsor/final\_paper/168\_hours/}. These prior simulations
demonstrate that hourly trial activity, daily AE escalation, and
weekly cumulative reporting all run on commodity hardware
\cite{kawchak_2026_19994945, kawchak_2026_19119939,
kawchak_2026_19396256}.

(e) The CLINICAL TREATMENT Act of 2020 (PL 116-260) covers routine
costs of qualifying trials for Medicaid beneficiaries from January 1,
2022, but does not establish a positive patient-initiated booking
right \cite{rank2_clinical_treatment_act_2020}.

(f) Affordable Care Act Section 2709 (PL 111-148, 2010) provides a
private-insurance baseline for routine clinical-trial cost coverage
but does not reach the decentralized-component payment-parity case
that HR 9501 closes \cite{rank3_aca_clinical_trials_2010,
usc_300gg8_clinical_trials}.

(g) Mazor 2025 (JAMA Network Open) reported a null result on AI
notifications alone in a randomized clinical trial of cancer
progression alerts; this finding informs the bill by indicating that
self-selection rights without a complete pathway from discovery to
consent to enrollment may not improve participation on their own
\cite{mazor2025ai_notifications_trial}.

\subsection{Definitions}

For the purposes of this Act, the following definitions apply.

(a) \billterm{Patient self-selection} in this Act means the patient's right to
discover, evaluate, and book the patient's own oncology clinical
trial slot without intermediary provider gatekeeping.\ The current operational
baseline is the existing 24/7 trial format documented in
\texttt{new-trial/national-24-7-trial/paper/full-paper/main.tex} and
the per-hour patient records files at
\texttt{new-trial/national-24-7-trial/hour-00/} through
\texttt{hour-55/} \cite{kawchak_2026_19994945}.

(b) \billterm{Machine-readable eligibility report} means a
structured FHIR-based output containing the patient's eligibility-relevant data fields for at least 80 percent of currently recruiting
United States oncology trials. The structured-eligibility baseline
is the ClinicalTrials.gov infrastructure under FDAAA Section 801
(PL 110-85, 2007) \cite{rank5_fdaaa_2007}.

(c) \billterm{Qualified Physical AI oncology trial site} in this specific Act designates a
site authorized under either in this context or under the prior California
SB 1042 / Title 22 framework regulation documented at
\texttt{new-trial/site/05-state-regulations/} and at
\texttt{national-platform/new\_trial\_psl/05\_title22\_regulations.tex}
\cite{kawchak_2026_19176370, kawchak_2026_19244918}.

(d) \billterm{24/7 booking} means the patient's right to book any
qualified site at any hour of the day, any day of the week. The
operational baseline is the 168-hour autonomous sponsor extension at
\texttt{sponsor/final\_paper/168\_hours/} \cite{kawchak_2026_19396256}.

(e) \billterm{AI matching tool} means an algorithm that proposes
trials given the patient's machine-readable eligibility report. The
qualified instances enumerated in this bill include TrialGPT, PRISM,
and TrialMatchAI \cite{jin2024trialgpt, gupta2024prism,
abdallah2026trialmatchai}; equivalent successors qualify under the
HTI-1 DSI transparency framework \cite{hti1_dsi_fact_sheet_2023}.

\subsection{Patient Self-Selection Rights}

\hspace{0.5cm} The patient holds the following rights under this Act, enforceable
against any CMS-billable entity, qualified Physical AI oncology trial
site, or sponsor.

(a) Right to receive a machine-readable eligibility report, in FHIR
format, from any CMS-billable EHR within seven calendar days of
pathology or genomic result release.

(b) Right to discover all currently recruiting United States oncology
trials via an FDA-published API that integrates ClinicalTrials.gov,
the FDA RTCT pilot directory, and the FDA Oncology Center of
Excellence Advancing Oncology Decentralized Trials directory
\cite{fda_oce_advancing_oncology_dct_2024, fda_rtct_press_2026}.

(c) Right to book a 24/7 continuous trial patient slot at any qualified physical AI oncology clinical trial site without
prior written referral.\ The operational pattern is documented in
\texttt{new-trial/national-24-7-trial/hour-NN/diagram\_facility.txt}
across hours 00, 12, 23, 47, and beyond \cite{kawchak_2026_19994945}.

(d) Right to payment parity for decentralized and at-home trial
components under 42 USC 300gg-8 and CMS NCD 310.1; this closes the
gap left by ACA Section 2709 \cite{usc_300gg8_clinical_trials,
cms_ncd_3101_routine_costs}.

(e) Right to cross-state portability of eligibility data and trial
bookings, drawing on the technical baseline in the ASTP/ONC Hospital
APIs Brief \cite{astp_onc_hospital_apis_2024}.

(f) Right to refuse provider-only routing of trial discovery in
favor of an AI matching tool of the patient's choosing (TrialGPT,
PRISM, TrialMatchAI), with disclosure of model provenance under
HTI-1 DSI transparency \cite{hti1_dsi_fact_sheet_2023}.

(g) Right to a single-click booking confirmation that creates an
audit-trail record at the clinical trial sponsor's pharmacovigilance endpoint, per
the architecture documented at
\texttt{sponsor/final\_paper/scripts/core\_agents/audit\_trail\_manager.py}
\cite{kawchak_2026_19396256}.

\subsection{Implementation}


\hspace{0.5cm} (a) The Secretary of Health and Human Services, acting through the
Food and Drug Administration, the Office of the National Coordinator
for Health Information Technology, and the Centers for Medicare and
Medicaid Services, shall publish the unified federal API described in
Subsection 1.3(b) within 12 months of the date of enactment.

(b) Each State, acting through its existing health-information
infrastructure, shall achieve full self-selection-portal coverage
across CMS-billable EHRs within 24 months of enactment.

(c) Implementation shall preserve stronger State-level patient
protections, including the California SB 37 routine-cost coverage
floor, under the preemption rule set forth in this Subsection
\cite{rank8_california_sb37_2001, rank4_fdora_fda_modernization_2022}.

(d) The FDA, ONC, HHS, and CMS shall jointly publish, within 18
months of enactment, an implementation guidance document that
references this Act and the FDA Decentralized Clinical Trials Final
Guidance 2024 \cite{fda_dct_guidance_2024}.

\subsection{Reporting and Enforcement}

\hspace{0.5cm} (a) The Secretary shall publish, annually, a report on patient self-selection rates by State, by payer, and by cancer type. The federal
Patient Control Dashboard's Domain 1 (Access) metrics anchor this
report: trial-offer rate, days from pathology or genomic result to
trial shortlist, machine-readable eligibility report coverage, real-time data release rates, and travel and work hours saved
\cite{hhs_right_to_access_2025, rocque2025remote_symptom_monitoring}.

(b) Information-blocking penalties under
\cite{federal_register_info_blocking_disincentives_2024} apply to any
CMS-billable entity that fails to expose the self-selection portal
endpoint required by Subsection 1.4.

(c) No patient shall be denied trial consideration by automated
output alone. Each automated rejection triggers a human reviewer
escalation under the HTI-1 DSI framework
\cite{rank9_onc_hti1_2023, hti1_dsi_fact_sheet_2023}.

\subsection{Patient-Control Framings}

\hspace{0.5cm}  For the individual patient, an NSCLC patient at home in rural Texas
can, on the patient's own initiative, receive an FHIR eligibility
report within seven days of the patient's genomic result, then
self-book a 24/7 slot at a USL-7-or-above surgical robot system, then
follow the 10-stage journey for PAT-2026-0042 documented at
\texttt{patient-journey/stage\_01\_prescreening.py} through
\texttt{stage\_10\_closeout.py} \cite{kawchak_2026_19119939}.

For broad adoption, all United States oncology trials must publish
eligibility criteria in machine-readable form within the 12-month
federal-API window, with all CMS-billable EHRs maintaining a self-selection portal endpoint accessible by every United States cancer
patient by month 24. The United States must remain Number 1 in the
world for cancer patient benefit during the medical AI and robotics
revolution. HR 9501 is the foundational bill that opens that
pathway, replacing the prior provider-gatekept trial referral
pattern with patient-initiated, AI-matched, 24/7-bookable self-
selection across all qualified Physical AI oncology trial sites
\cite{kawchak_2026_19994945, kawchak_2026_19119939,
rank2_clinical_treatment_act_2020}.
